\begin{explanationbox}
	\textbf{\textcolor{CExplanation}{一句话直觉}}
	平移异面直线使其相交，构造三角形后，用余弦定理求对应角余弦，取绝对值确保所成角为锐角或直角。

	\medskip
	\textbf{\textcolor{CExplanation}{核心拆解}}
	\begin{description}[style=nextline, leftmargin=2.8em, labelsep=0.5em, itemsep=0.45em]
		\item[\textbf{要点一（平移共起点）：}]
			将一条异面直线平移，使其与另一条相交于一点，两线段从同一端点出发。
		\item[\textbf{要点二（余弦定理）：}]
			在平移后得到的三角形中，已知三边长度，余弦定理直接给出任一内角的余弦值。
		\item[\textbf{要点三（绝对值取锐角）：}]
			异面直线所成角范围是$(0^\circ,90^\circ]$，余弦非负；若内角为钝角，其余弦为负，取绝对值后得到锐角余弦。
	\end{description}

	\medskip
	\textbf{\textcolor{CExplanation}{几何本质（空间转平面）}}
	通过平移将异面直线转化为共面相交直线，将空间角度测量问题直接转化为平面三角形角度计算问题。

	\medskip
	\textbf{\textcolor{CExplanation}{代数意义（余弦变式）}}
	公式 $\cos\theta = \dfrac{|m^2+n^2-d^2|}{2mn}$ 是余弦定理的变形，分子为两边平方和减去第三边平方的绝对值，表示两线段长度关系与夹角余弦的联系。

	\medskip
	\textbf{\textcolor{CExplanation}{考点价值（直接与逆向）}}
	\begin{itemize}[leftmargin=*, itemsep=0.5em, topsep=0.4em]
		\item \textbf{考法一（直接套公式）：} 在给定几何体中，通过平移找出三角形三边，直接代入公式求角。
		\item \textbf{考法二（逆向应用）：} 已知角与两边，求第三边；或建立等量关系证明角度相等。
	\end{itemize}

	\medskip
	\textbf{\textcolor{CExplanation}{顿悟点}}
	\textbf{平移是关键，余弦定理是工具，绝对值是保障。}

	\medskip
	\textbf{\textcolor{CExplanation}{使用场景}}
	\begin{itemize}[leftmargin=*, itemsep=0.5em, topsep=0.4em]
		\item \textbf{场景一（三边已知）：} 在几何体中，能轻易找到平移后的三角形三边长度，直接求角。
		\item \textbf{场景二（已知角与两边）：} 逆用公式，已知所求角与两线段长度，求第三边或验证数据。
	\end{itemize}

\end{explanationbox}
